\chapter{几何组成分析}

\section{判别方法速记}

\begin{notecard}
几何组成分析的核心不是套公式，而是判断体系能否作为刚片逐步并入大地。常用顺序如下：
\begin{enumerate}
  \item 先识别大地、刚片、链杆、铰和支座约束。
  \item 对平面体系，可先用计算自由度
  \[
    W=3m-2h-r
  \]
  作必要检查，其中 \(m\) 为刚片数，\(h\) 为单铰数，\(r\) 为支座链杆数或外部约束数。
  \item 再用二元体规则、两刚片规则和三刚片规则作充分性判断。
  \item 若约束方向平行且等长，常见为常变；平行但不等长，常见为瞬变。若已经几何不变但约束超过必要数，则存在多余约束。
\end{enumerate}
\end{notecard}

\section{第 1--6 题}

\begin{problemcard}{1-1 几何组成判别 \source{浙江大学 2008}}
图 1-56 所示体系的几何组成为（\quad）。

\begin{tabularx}{\textwidth}{@{}XX@{}}
A. 几何瞬变，无多余约束 & B. 几何不变 \\
C. 几何常变 & D. 几何瞬变，有多余约束
\end{tabularx}

\begin{figure}[H]
\centering
\begin{tikzpicture}[scale=0.92]
  \draw[rigid] (0,2) rectangle (3.9,2.35);
  \draw[rigid] (0,0) rectangle (3.9,0.35);
  \draw[rigid] (0,0.35) rectangle (0.38,2);
  \foreach \x/\dx in {0.75/0,1.55/0.08,2.45/0.22,3.65/0.38}{
    \node[pin] at (\x,2) {};
    \node[pin] at (\x+\dx,0.35) {};
    \draw[link] (\x,2)--(\x+\dx,0.35);
  }
  \figtitle{图 1-56：四根链杆近似平行且不等长}
\end{tikzpicture}
\end{figure}
\end{problemcard}

\begin{solutioncard}
\textbf{答案：}\answer{D}

\textbf{解析：}把上下两个阴影部分视为两块刚片。两刚片之间由 4 根链杆联系，若要形成无多余约束的几何不变体系，平面两刚片之间通常需要 3 个相互独立的约束；这里有 4 根链杆，因此从数量上看存在多余约束。

进一步看约束方向：4 根链杆方向近似平行，但长度不等，不能形成稳定的几何约束组合。按讲解笔记中的口诀，“平行且等长为常变，平行不等长为瞬变”，故该体系为几何瞬变，并且有两个多余约束。
\end{solutioncard}

\begin{problemcard}{1-2 补加支座约束 \source{中南大学 2011；合肥工业大学 2009}}
欲使图 1-57 所示体系成为无多余约束的几何不变体系，则需在 \(A\) 端加入（\quad）。

\begin{tabularx}{\textwidth}{@{}XXXX@{}}
A. 固定铰支座 & B. 固定支座 & C. 滑动铰支座 & D. 定向支座
\end{tabularx}

\begin{figure}[H]
\centering
\begin{tikzpicture}[scale=0.88]
  \coordinate (A) at (2.9,1.15);
  \coordinate (B) at (2.35,2.4);
  \coordinate (C) at (3.45,2.4);
  \coordinate (D) at (2.35,3.05);
  \coordinate (F) at (3.45,3.05);
  \coordinate (E) at (0.65,0.25);
  \coordinate (G) at (5.15,0.25);
  \draw[member] (0.6,0.25)--(0.6,3.05)--(D);
  \draw[member] (D)--(B)--(C)--(F)--(5.2,3.05)--(5.2,0.25);
  \draw[member] (B)--(A)--(C);
  \draw[member] (A)--(2.9,0.75);
  \node[pin] at (B) {};
  \node[pin] at (C) {};
  \node[pin] at (D) {};
  \node[pin] at (F) {};
  \node[pin,label=below left:\(E\)] at (E) {};
  \node[pin,label=below right:\(G\)] at (G) {};
  \roller{0.6}{0.25}
  \roller{5.2}{0.25}
  \fixedbase{2.9}{0.75}
  \node[below right] at (A) {\(A\)};
  \node[above left] at (B) {\(B\)};
  \node[above right] at (C) {\(C\)};
  \node[above] at (D) {\(D\)};
  \node[above] at (F) {\(F\)};
  \figtitle{图 1-57：在 \(A\) 处补固定端}
\end{tikzpicture}
\end{figure}
\end{problemcard}

\begin{solutioncard}
\textbf{答案：}\answer{B}

\textbf{解析：}先作自由度必要检查。原体系可取 3 个刚片、1 个单铰和 4 个外部约束，则
\[
  W=3\times 3-2\times 1-4=3 .
\]
要使体系成为几何不变体系，还需要补足 3 个独立约束。固定端正好提供两个平动约束和一个转动约束，共 3 个约束，因此候选项为固定支座。

加入固定端后，刚片 \(ABC\) 先与大地合并；随后刚片 \(DE\) 可通过链杆 \(BD\) 与铰 \(E\) 并入大地，刚片 \(FG\) 亦可用同样方式并入大地。全过程没有额外约束，故为无多余约束的几何不变体系。
\end{solutioncard}

\begin{problemcard}{1-3 两刚片规则 \source{中南大学 2009}}
图 1-58 所示体系的几何构造性质是（\quad）。

\begin{tabularx}{\textwidth}{@{}XX@{}}
A. 常变体系 & B. 瞬变体系 \\
C. 无多余联系的几何不变体系 & D. 有多余联系的几何不变体系
\end{tabularx}

\begin{figure}[H]
\centering
\begin{tikzpicture}[scale=0.78]
  \coordinate (A) at (3.1,0.7);
  \coordinate (B) at (1.55,3.0);
  \coordinate (C) at (4.75,3.0);
  \coordinate (D) at (0.55,0.35);
  \coordinate (E) at (5.75,0.35);
  \draw[member] (0.5,0.35)--(0.5,2.0)--(B)--(2.75,1.75)--(2.75,0.7);
  \draw[member] (2.75,1.75)--(C)--(5.8,2.0)--(5.8,0.35);
  \draw[link,dashed] (D)--(B);
  \draw[link,dashed] (E)--(C);
  \node[pin,label=below:\(A\)] at (A) {};
  \wallroller{3.1}{0.7}
  \node[pin,label=above:\(B\)] at (B) {};
  \node[pin,label=above:\(C\)] at (C) {};
  \node[pin,label=below:\(D\)] at (D) {};
  \node[pin,label=below:\(E\)] at (E) {};
  \roller{0.5}{0.35}
  \roller{5.8}{0.35}
  \figtitle{图 1-58：三根不共线且不交于一点的约束}
\end{tikzpicture}
\end{figure}
\end{problemcard}

\begin{solutioncard}
\textbf{答案：}\answer{C}

\textbf{解析：}将上部折线框架视为刚片 \(ABC\)，大地视为另一刚片。二者之间由三根链杆联系：\(A\)、\(BD\)、\(CE\)。这三根链杆不共线，且其作用线不交于同一点，满足两刚片规则。

由于正好提供使两刚片相对固定所需的三个独立约束，没有多余联系，因此该体系为无多余联系的几何不变体系。
\end{solutioncard}

\begin{problemcard}{1-4 常变体系与多余约束 \source{广西大学 2010}}
图 1-59 所示体系内部是 \underline{\hspace{3em}} 体系，它有 \underline{\hspace{3em}} 个多余约束。

\begin{figure}[H]
\centering
\begin{tikzpicture}[scale=0.8]
  \foreach \x in {0,1.3,2.6,3.9,5.2}{
    \node[pin] at (\x,2) {};
  }
  \foreach \x in {0,1.3,2.6,3.9,5.2}{
    \node[pin] at (\x,0) {};
  }
  \draw[member] (0,2)--(5.2,2)--(5.2,0)--(0,0)--cycle;
  \draw[member] (1.3,2)--(1.3,0);
  \draw[member] (2.6,2)--(2.6,0);
  \draw[member] (3.9,2)--(3.9,0);
  \draw[member] (0,2)--(1.3,0);
  \draw[member] (1.3,2)--(2.6,0);
  \draw[member] (3.9,2)--(5.2,0);
  \draw[member] (3.9,0)--(5.2,2);
  \wallroller{0}{0}
  \roller{5.2}{0}
  \figtitle{图 1-59：内部桁式体系}
\end{tikzpicture}
\end{figure}
\end{problemcard}

\begin{solutioncard}
\textbf{答案：}\answer{几何常变；1}

\textbf{解析：}采用两刚片规则分析。可把左侧部分看成刚片 \(ABCD\)，右侧部分看成刚片 \(EFGH\)。右侧刚片内部本身已有一个多余约束，而左右两刚片之间仅由链杆 \(CE\) 和 \(DF\) 联系。

两刚片若要组成几何不变整体，需要三根不共线且不交于一点的链杆，或等效的三个独立约束。这里只有两根链杆，不能限制两个刚片之间的全部相对运动，因此体系为几何常变。由于右侧刚片内部已有一个多余约束，故多余约束数为 1。
\end{solutioncard}

\begin{problemcard}{1-5 计算自由度 \source{哈尔滨工业大学 2011}}
图 1-60 所示体系的计算自由度 \(W=\underline{\hspace{4em}}\)。

\begin{figure}[H]
\centering
\begin{tikzpicture}[scale=0.86]
  \coordinate (P0) at (0,1.25);
  \coordinate (P1) at (1.6,1.25);
  \coordinate (P2) at (3.2,1.25);
  \coordinate (P3) at (4.6,1.25);
  \coordinate (P4) at (4.6,3.0);
  \coordinate (P5) at (6.0,1.25);
  \draw[member] (P0)--(P1)--(P2)--(P3)--(P5)--(P4)--(P3);
  \draw[member] (P2)--(3.2,0.05);
  \draw[member] (P2)--(4.6,0.05);
  \draw[member] (P4)--(P5);
  \node[pin] at (P0) {};
  \node[pin] at (P1) {};
  \node[pin] at (P2) {};
  \node[pin] at (P3) {};
  \node[pin] at (P4) {};
  \node[pin] at (P5) {};
  \wallroller{0}{1.25}
  \fixedbase{3.2}{0.05}
  \roller{4.6}{0.05}
  \node[above] at (3.2,1.35) {5};
  \node[above] at (4.05,1.35) {2};
  \node[right] at (3.75,0.65) {3};
  \node[left] at (3.2,0.65) {4};
  \node[right] at (5.2,2.0) {1};
  \figtitle{图 1-60：按刚片编号计算自由度}
\end{tikzpicture}
\end{figure}
\end{problemcard}

\begin{solutioncard}
\textbf{答案：}\answer{\(-4\)}

\textbf{解析：}按答案图所示将体系划分为 5 个刚片，体系中共有 6 个单铰，支座约束折算为 4 个外部约束，并有 1 个复铰折算为 3 个约束。于是计算自由度为
\[
  W = 5\times 3 - 6\times 2 - 1\times 3 - 4 = -4 .
\]
计算自由度为负说明约束数超过自由度数，体系中存在多余约束。注意：\(W<0\) 只是必要检查结果，最终几何组成仍需结合刚片规则判断。
\end{solutioncard}

\begin{problemcard}{1-6 二元体规则 \source{中南大学 2012}}
图 1-61 所示体系几何不变且有多余联系。（\quad）

\begin{figure}[H]
\centering
\begin{tikzpicture}[scale=0.74]
  \coordinate (A) at (0,0); \coordinate (B) at (1.4,0);
  \coordinate (C) at (1.4,1.25); \coordinate (D) at (0,1.25);
  \coordinate (E) at (1.4,2.5); \coordinate (F) at (0,2.5);
  \coordinate (G) at (2.9,2.5); \coordinate (H) at (2.9,0);
  \coordinate (I) at (1.4,3.75); \coordinate (J) at (0,3.75);
  \coordinate (K) at (-1.2,3.1); \coordinate (L) at (-1.2,2.5);
  \coordinate (M) at (-2.5,2.5);
  \draw[member] (A)--(B)--(C)--(D)--cycle;
  \draw[member] (D)--(F)--(E)--(C);
  \draw[member] (C)--(E);
  \draw[member] (F)--(J)--(I)--(E);
  \draw[member] (E)--(G)--(H);
  \draw[member] (I)--(G);
  \draw[member] (M)--(L)--(F);
  \draw[member] (M)--(K)--(J);
  \draw[member] (K)--(L);
  \draw[member] (K)--(F);
  \foreach \p/\lab in {A/A,B/B,C/C,D/D,E/E,F/F,G/G,H/H,I/I,J/J,K/K,L/L,M/M}{
    \node[pin,label=above:{\scriptsize \(\lab\)}] at (\p) {};
  }
  \wallroller{0}{0}
  \roller{0}{0}
  \roller{1.4}{0}
  \roller{2.9}{0}
  \figtitle{图 1-61：可由大地逐步添加二元体}
\end{tikzpicture}
\end{figure}
\end{problemcard}

\begin{solutioncard}
\textbf{答案：}\answer{错误，记作 \(\times\)}

\textbf{解析：}该体系可由大地开始，依次添加二元体
\[
  AB,\ ACB,\ ADC,\ DEC,\ DFE,\ EGH,\ EIG,\ FJI,\ FKJ,\ KLF,\ KML .
\]
每添加一个二元体，都是用两根不共线链杆把一个新节点稳定地接入已有几何不变部分，因此不会引入多余约束。逐步添加完成后，整体为无多余约束的几何不变体系。

题干说“几何不变且有多余联系”，前半句正确，后半句错误，所以判断题答案为 \(\times\)。
\end{solutioncard}

\section{第 7--13 题}

\begin{problemcard}{1-7 刚片并入与多余约束 \source{哈尔滨工业大学 2009}}
对图 1-62 所示体系进行几何组成分析。

\begin{figure}[H]
\centering
\begin{tikzpicture}[scale=0.9]
  \coordinate (A) at (0,0); \coordinate (B) at (1.55,0);
  \coordinate (C) at (3.1,0); \coordinate (D) at (0,2);
  \coordinate (E) at (1.55,2); \coordinate (F) at (3.1,2);
  \draw[member] (A)--(D)--(E)--(F)--(C);
  \draw[member] (B)--(E);
  \foreach \p/\lab in {A/A,B/B,C/C,D/D,E/E,F/F}{
    \node[pin,label=above:{\scriptsize \(\lab\)}] at (\p) {};
  }
  \fixedbase{0}{0}
  \fixedbase{1.55}{0}
  \fixedbase{3.1}{0}
  \figtitle{图 1-62：三立柱框架体系}
\end{tikzpicture}
\end{figure}
\end{problemcard}

\begin{solutioncard}
\textbf{结论：}\answer{有 3 个多余约束的几何不变体系}

\textbf{解析：}刚片 \(AD\)、\(BE\) 可直接并入大地，链杆 \(DE\) 因而成为一个多余约束。随后刚片 \(EFC\) 与大地通过固定端 \(C\) 和铰 \(E\) 联系，满足两刚片规则；由于固定端已提供较强约束，此处还会带来两个多余约束。

因此整体虽然几何不变，但共有 \(1+2=3\) 个多余约束。
\end{solutioncard}

\begin{problemcard}{1-8 三刚片规则 \source{武汉科技大学 2009}}
对图 1-63 所示体系作几何组成分析。

\begin{figure}[H]
\centering
\begin{tikzpicture}[scale=0.72]
  \coordinate (A) at (0,0); \coordinate (B) at (0,2.7);
  \coordinate (C) at (2.2,2.7); \coordinate (D) at (4.4,2.7);
  \coordinate (E) at (4.4,0); \coordinate (G) at (2.2,0);
  \coordinate (F) at (2.2,1.55);
  \draw[member] (A)--(B)--(C)--(D)--(E)--(G)--(A);
  \draw[member] (B)--(F)--(D);
  \draw[member] (A)--(C)--(E);
  \draw[member] (B)--(G)--(D);
  \draw[member] (C)--(F)--(G);
  \draw[member,dashed] (-1.0,-0.35)--(F)--(5.6,-0.35);
  \node[left] at (-1.0,-0.35) {\scriptsize \(O_1\)};
  \node[right] at (5.6,-0.35) {\scriptsize \(O_2\)};
  \foreach \p/\lab in {A/A,B/B,C/C,D/D,E/E,F/F,G/G}{
    \node[pin,label=above:{\scriptsize \(\lab\)}] at (\p) {};
  }
  \wallroller{0}{0}
  \roller{0}{0}
  \roller{4.4}{0}
  \figtitle{图 1-63：三铰不共线的三刚片体系}
\end{tikzpicture}
\end{figure}
\end{problemcard}

\begin{solutioncard}
\textbf{结论：}\answer{无多余约束的几何不变体系}

\textbf{解析：}按三刚片规则分析。刚片 \(ABC\) 与刚片 \(CDE\) 交于铰 \(C\)；刚片 \(ABC\) 与刚片 \(FG\) 由链杆 \(BF\)、\(AG\) 联系，其作用线交于虚铰 \(O_2\)；刚片 \(CDE\) 与刚片 \(FG\) 由链杆 \(DF\)、\(GE\) 联系，其作用线交于虚铰 \(O_1\)。

三个铰 \(C,O_1,O_2\) 不共线，满足三刚片规则，可组成一个大刚片；该大刚片再由三根不共线且不交于一点的链杆支座与大地联系，满足两刚片规则，所以体系几何不变且无多余约束。
\end{solutioncard}

\begin{problemcard}{1-9 三铰共线瞬变 \source{武汉理工大学 2008}}
对图 1-64 所示杆件体系进行几何构造分析。

\begin{figure}[H]
\centering
\begin{tikzpicture}[scale=0.78]
  \coordinate (A) at (0,2); \coordinate (B) at (1.3,2);
  \coordinate (C) at (1.3,0.85); \coordinate (D) at (4.1,2);
  \coordinate (E) at (5.4,2); \coordinate (F) at (4.1,0.85);
  \coordinate (G) at (2.7,0);
  \draw[member] (A)--(B)--(D)--(E);
  \draw[member] (B)--(C)--(F)--(D);
  \draw[member] (A)--(C);
  \draw[member] (D)--(F)--(G)--(C);
  \draw[member] (A)--(G)--(E);
  \draw[member,dashed] (E)--(6.4,2);
  \draw[member,dashed] (F)--(6.4,0.85);
  \node[right] at (6.4,2) {\scriptsize \(O_1\)};
  \foreach \p/\lab in {A/A,B/B,C/C,D/D,E/E,F/F,G/G}{
    \node[pin,label=above:{\scriptsize \(\lab\)}] at (\p) {};
  }
  \roller{0}{2}
  \roller{5.4}{2}
  \roller{2.7}{0}
  \figtitle{图 1-64：三铰共线导致瞬变}
\end{tikzpicture}
\end{figure}
\end{problemcard}

\begin{solutioncard}
\textbf{结论：}\answer{有 1 个多余约束的几何瞬变体系}

\textbf{解析：}取刚片 \(ABC\)、刚片 \(DEF\) 和大地为三刚片。刚片 \(ABC\) 与刚片 \(DEF\) 由链杆 \(BD\)、\(CF\) 联系，二者作用线平行，可视为交于无穷远铰 \(O_1\)；刚片 \(ABC\) 与大地由链杆 \(A\)、\(CG\) 联系，交于铰 \(A\)；刚片 \(DEF\) 与大地由链杆 \(FG\)、\(E\) 联系，交于铰 \(E\)。

三个等效铰共线，不满足三刚片规则，所以体系瞬变。由于约束数量超过必要数，体系同时有 1 个多余约束。
\end{solutioncard}

\begin{problemcard}{1-10 计算自由度并判别 \source{华南理工大学 2012}}
计算图 1-65 所示体系的计算自由度，并分析其几何组成。

\begin{figure}[H]
\centering
\begin{tikzpicture}[scale=0.73]
  \foreach \x in {0,1.2,2.4,3.6,4.8}{
    \node[pin] at (\x,2.2) {};
    \node[pin] at (\x,0) {};
  }
  \node[pin] at (1.2,1.1) {};
  \node[pin] at (3.6,1.1) {};
  \draw[member] (0,0)--(4.8,0)--(4.8,2.2)--(0,2.2)--cycle;
  \draw[member] (1.2,2.2)--(1.2,1.1)--(0,0);
  \draw[member] (1.2,1.1)--(2.4,2.2)--(3.6,1.1)--(4.8,2.2);
  \draw[member] (1.2,1.1)--(2.4,0)--(3.6,1.1)--(4.8,0);
  \draw[member] (0,2.2)--(1.2,1.1);
  \draw[member] (3.6,2.2)--(3.6,1.1);
  \wallroller{0}{0}
  \roller{0}{0}
  \roller{4.8}{0}
  \figtitle{图 1-65：有一个多余约束的几何不变体系}
\end{tikzpicture}
\end{figure}
\end{problemcard}

\begin{solutioncard}
\textbf{结论：}\answer{\(W=-1\)，有 1 个多余约束的几何不变体系}

\textbf{解析：}按题图统计，计算自由度为
\[
  W=12\times 2-25=-1 .
\]
再作几何组成判断：刚片 \(ABCD\) 与刚片 \(DEFB\) 通过铰 \(B\) 和铰 \(D\) 联系，相当于四个约束。两刚片规则只需三个独立约束即可使二者相对固定，因此满足几何不变条件，同时多出一个约束。
\end{solutioncard}

\begin{problemcard}{1-11 超静定次数 \source{青岛理工大学 2012}}
对图 1-66 所示平面体系进行几何组成分析，并指出结构超静定次数。

\begin{figure}[H]
\centering
\begin{tikzpicture}[scale=0.78]
  \coordinate (A) at (0,0); \coordinate (B) at (0,2.2);
  \coordinate (C) at (2.4,2.2); \coordinate (D) at (2.4,0);
  \coordinate (E) at (4.8,0); \coordinate (F) at (4.8,2.2);
  \coordinate (G) at (6.0,2.2); \coordinate (H) at (6.0,0);
  \draw[member] (A)--(B)--(C)--(F)--(G);
  \draw[member] (C)--(D);
  \draw[member] (E)--(F);
  \draw[member] (D)--(E)--(H);
  \draw[member] (B)--(D);
  \draw[member] (D)--(C);
  \draw[member] (C)--(E);
  \draw[member] (F)--(H);
  \foreach \p/\lab in {A/A,B/B,C/C,D/D,E/E,F/F,G/G,H/H}{
    \node[pin,label=above:{\scriptsize \(\lab\)}] at (\p) {};
  }
  \wallroller{0}{0}
  \roller{0}{0}
  \roller{2.4}{0}
  \roller{6.0}{0}
  \wallroller{6.0}{2.2}
  \figtitle{图 1-66：两次超静定体系}
\end{tikzpicture}
\end{figure}
\end{problemcard}

\begin{solutioncard}
\textbf{结论：}\answer{有 2 个多余约束的几何不变体系，即两次超静定}

\textbf{解析：}先去掉二元体 \(AB\) 和 \(CEG\)，不改变体系基本几何性质。剩余部分可看作带有两个多余约束的刚片 \(BCD\)，它与大地由三根不共线且不交于一点的链杆 \(D\)、\(CF\)、\(CG\) 联系，满足两刚片规则。

因此体系几何不变，同时存在两个多余约束，结构超静定次数为 2。
\end{solutioncard}

\begin{problemcard}{1-12 几何稳定性分析 \source{东南大学 2008}}
分析图 1-67 所示体系的几何稳定性，写出分析过程。

\begin{figure}[H]
\centering
\begin{tikzpicture}[scale=0.82]
  \coordinate (A) at (0,0); \coordinate (B) at (3.2,0);
  \coordinate (F) at (6.4,0); \coordinate (C) at (1.5,1.55);
  \coordinate (E) at (4.9,1.55); \coordinate (D) at (0.6,3.0);
  \coordinate (G) at (5.8,3.0);
  \draw[member] (A)--(B)--(F);
  \draw[member] (A)--(C)--(B)--(E)--(F);
  \draw[member] (C)--(D);
  \draw[member] (E)--(G);
  \draw[member] (C)--(1.5,0);
  \draw[member] (E)--(4.9,0);
  \foreach \x in {1.0,2.1,4.2,5.3}{
    \node[smallpin] at (\x,0.35) {};
    \draw[member] (\x,0.35)--(\x,0);
  }
  \foreach \p/\lab in {A/A,B/B,C/C,D/D,E/E,F/F,G/G}{
    \node[pin,label=above:{\scriptsize \(\lab\)}] at (\p) {};
  }
  \wallroller{0}{0}
  \roller{0}{0}
  \roller{6.4}{0}
  \wallroller{0.6}{3.0}
  \wallroller{5.8}{3.0}
  \figtitle{图 1-67：含一根多余杆 \(GE\)}
\end{tikzpicture}
\end{figure}
\end{problemcard}

\begin{solutioncard}
\textbf{结论：}\answer{有 1 个多余约束的几何不变体系}

\textbf{解析：}刚片 \(ABC\) 与大地由三根不共线且不交于一点的链杆支座 \(A\) 和链杆 \(CD\) 联系，满足两刚片规则，可先并入大地。随后刚片 \(FBE\) 与大地由铰 \(B\) 和链杆 \(F\) 联系，也满足两刚片规则。

在上述并入过程完成后，杆 \(GE\) 已不再是维持几何不变所必需的约束，因此为多余约束。体系为有一个多余约束的几何不变体系。
\end{solutioncard}

\begin{problemcard}{1-13 等效刚片分析 \source{山东建筑大学 2011}}
试对图 1-68 所示体系作几何组成分析。

\begin{figure}[H]
\centering
\begin{tikzpicture}[scale=0.88]
  \coordinate (P1) at (0,0);
  \coordinate (P2) at (1.3,0);
  \coordinate (P4) at (2.8,0);
  \coordinate (P5) at (3.6,-1.0);
  \coordinate (P6) at (4.8,-1.0);
  \coordinate (P3) at (5.8,0);
  \draw[member] (P1)--(P2)--(P4)--(P3);
  \draw[member] (P4)--(P5)--(P6)--(P3);
  \draw[member] (P4)--(P6);
  \draw[member] (P5)--(P3);
  \foreach \p/\lab in {P1/1,P2/2,P3/3,P4/4,P5/5,P6/6}{
    \node[pin,label=above:{\scriptsize \(\lab\)}] at (\p) {};
  }
  \roller{0}{0}
  \roller{1.3}{0}
  \roller{5.8}{0}
  \wallroller{5.8}{0}
  \node[font=\scriptsize,Muted] at (4.35,-1.55) {刚片 \(3456\) 可等效为链杆 \(34\)};
  \figtitle{图 1-68：等效为两刚片规则判断}
\end{tikzpicture}
\end{figure}
\end{problemcard}

\begin{solutioncard}
\textbf{结论：}\answer{无多余约束的几何不变体系}

\textbf{解析：}刚片 \(3456\) 可等效为链杆 \(34\)。等效后，取刚片 \(14\) 与大地分析：二者由三根不共线且不交于一点的链杆联系，满足两刚片规则。

约束数正好满足几何不变所需的独立约束数，因此体系为无多余约束的几何不变体系。
\end{solutioncard}

\section{第 14--17 题}

\begin{problemcard}{1-14 二元体去除法 \source{山东建筑大学 2008}}
试对图 1-69 所示体系作几何组成分析。

\begin{figure}[H]
\centering
\begin{tikzpicture}[scale=0.82]
  \coordinate (A) at (0,0); \coordinate (B) at (1.45,0);
  \coordinate (C) at (2.9,0); \coordinate (D) at (4.35,0);
  \coordinate (E) at (5.8,0);
  \coordinate (J) at (0,2.0); \coordinate (K) at (1.45,2.0);
  \coordinate (L) at (2.9,2.0); \coordinate (M) at (4.35,2.0);
  \coordinate (N) at (5.8,2.0);
  \draw[member] (A)--(J)--(K)--(B)--(A);
  \draw[member] (K)--(L)--(C)--(B);
  \draw[member] (L)--(M)--(D)--(C);
  \draw[member] (M)--(N)--(E)--(D);
  \draw[member] (J)--(B);
  \draw[member] (K)--(C);
  \draw[member] (M)--(C);
  \foreach \p in {A,B,C,D,E,J,K,L,M,N}{\node[pin] at (\p) {};}
  \wallroller{0}{0}
  \roller{0}{0}
  \roller{1.45}{0}
  \roller{2.9}{0}
  \wallroller{4.35}{0}
  \roller{5.8}{0}
  \figtitle{图 1-69：可先去除若干二元体}
\end{tikzpicture}
\end{figure}
\end{problemcard}

\begin{solutioncard}
\textbf{结论：}\answer{无多余约束的几何不变体系}

\textbf{解析：}先去掉二元体 \(ABC\)、\(AJ\)、\(DEF\) 和 \(GFH\) 等不改变几何性质的局部单元。剩余主刚片 \(JCDI\) 与大地由三根不共线且不交于一点的链杆支座联系，满足两刚片规则。

整个体系可由大地逐步并入，且没有超过必要数的独立约束，因此为无多余约束的几何不变体系。
\end{solutioncard}

\begin{problemcard}{1-15 计算自由度与常变体系 \source{中国地震局工程力学研究所 2010}}
计算图 1-70 的计算自由度，并作几何构造分析。

\begin{figure}[H]
\centering
\begin{tikzpicture}[scale=0.88]
  \coordinate (A) at (0,0); \coordinate (B) at (1.1,0);
  \coordinate (C) at (2.2,0); \coordinate (D) at (3.3,0);
  \coordinate (E) at (4.4,0);
  \coordinate (P) at (2.2,1.5);
  \coordinate (Q) at (3.3,1.5);
  \draw[member] (A)--(B)--(P)--(Q)--(D)--(E);
  \draw[member] (B)--(C)--(D);
  \draw[member] (B)--(Q);
  \draw[member] (P)--(D);
  \draw[member] (P)--(C);
  \draw[member] (Q)--(C);
  \foreach \p in {A,B,C,D,E,P,Q}{\node[pin] at (\p) {};}
  \roller{0}{0}
  \roller{4.4}{0}
  \figtitle{图 1-70：自由度为正的常变体系}
\end{tikzpicture}
\end{figure}
\end{problemcard}

\begin{solutioncard}
\textbf{结论：}\answer{\(W=1\)，无多余约束的几何常变体系}

\textbf{解析：}按题图统计可得计算自由度
\[
  W=7\times 3-2\times 10=1 .
\]
自由度为正，说明体系仍有独立运动可能。

再用两刚片规则分析：刚片 \(ABC\) 与刚片 \(DEF\) 仅由链杆 \(CD\)、\(BE\) 联系，独立约束不足三个，不满足两刚片规则。因此体系不是几何不变体系，而是无多余约束的几何常变体系。
\end{solutioncard}

\begin{problemcard}{1-16 三刚片共线判别 \source{西南交通大学 2008}}
对图 1-71 所示体系进行几何构造分析。

\begin{figure}[H]
\centering
\begin{tikzpicture}[scale=0.82]
  \coordinate (A) at (0,0); \coordinate (B) at (1.6,0);
  \coordinate (C) at (3.2,0); \coordinate (D) at (4.8,0);
  \coordinate (E) at (1.6,1.3); \coordinate (F) at (3.2,1.3);
  \coordinate (G) at (2.4,2.65); \coordinate (H) at (2.4,-1.25);
  \draw[member] (A)--(B)--(C)--(D);
  \draw[member] (A)--(E)--(F)--(D);
  \draw[member] (B)--(E);
  \draw[member] (C)--(F);
  \draw[member] (E)--(G)--(F);
  \draw[member] (B)--(H)--(C);
  \draw[member] (E)--(C);
  \draw[member] (B)--(F);
  \foreach \p/\lab in {A/A,B/B,C/C,D/D,E/E,F/F,G/G,H/H}{
    \node[pin,label=above:{\scriptsize \(\lab\)}] at (\p) {};
  }
  \roller{0}{0}
  \roller{4.8}{0}
  \roller{2.4}{-1.25}
  \wallroller{4.8}{0}
  \figtitle{图 1-71：三等效铰共线}
\end{tikzpicture}
\end{figure}
\end{problemcard}

\begin{solutioncard}
\textbf{结论：}\answer{有 1 个多余约束的几何瞬变体系}

\textbf{解析：}先去掉二元体 \(ABC\)。随后按三刚片规则分析：刚片 \(ADE\) 与大地由链杆 \(D\)、\(EH\) 联系，作用线交于无穷远铰 \(O_1\)；刚片 \(CFG\) 与大地由链杆 \(F\)、\(GH\) 联系，交于铰 \(O_2\)；刚片 \(ADE\) 与刚片 \(CFG\) 由链杆 \(AC\)、\(EG\) 联系，交于无穷远铰 \(O_3\)。

三个等效铰共线，不满足三刚片规则，体系瞬变；同时约束数量超过必要数，有 1 个多余约束。
\end{solutioncard}

\begin{problemcard}{1-17 计算自由度与多余约束 \source{北京工业大学 2010}}
计算图 1-72 所示体系的计算自由度，并进行几何组成分析。

\begin{figure}[H]
\centering
\begin{tikzpicture}[scale=0.85]
  \coordinate (A) at (0,3.0); \coordinate (B) at (2.3,3.0);
  \coordinate (C) at (2.3,1.55); \coordinate (D) at (4.4,1.55);
  \coordinate (E) at (2.3,0); \coordinate (P) at (0,0);
  \draw[member] (A)--(B)--(C)--(E);
  \draw[member] (C)--(D);
  \draw[member] (E)--(P);
  \foreach \p/\lab in {A/A,B/B,C/C,D/D,E/E}{
    \node[pin,label=above:{\scriptsize \(\lab\)}] at (\p) {};
  }
  \wallroller{0}{3.0}
  \roller{2.3}{0}
  \wallroller{0}{0}
  \wallroller{4.4}{1.55}
  \node[font=\scriptsize,text=Muted,anchor=west] at (4.65,2.05) {杆 \(CE\) 为多余约束};
  \figtitle{图 1-72：有一个多余约束的几何不变体系}
\end{tikzpicture}
\end{figure}
\end{problemcard}

\begin{solutioncard}
\textbf{结论：}\answer{\(W=-1\)，有 1 个多余约束的几何不变体系}

\textbf{解析：}按刚片和约束统计，计算自由度为
\[
  W=4\times 3-3\times 2-7\times 1=-1 .
\]
先将刚片 \(AB\) 并入大地；随后杆 \(CD\) 与大地由三根不共线且不交于一点的链杆 \(BC\) 和支座 \(D\) 联系，满足两刚片规则。杆 \(CE\) 在此基础上不再是维持几何不变所必需的约束，因此为多余约束。

所以体系为有一个多余约束的几何不变体系。
\end{solutioncard}

